\documentclass[10pt,a4paper]{amsart}
\usepackage[margin=19mm]{geometry}
\usepackage{amsmath,amssymb,mathtools}
\usepackage[scr=rsfs,cal=euler]{mathalfa}
\usepackage{xfrac}
\usepackage[unicode,hidelinks,bookmarks=false]{hyperref}
\newcommand{\ie}{i.e.\ }
\newcommand{\cf}{cf.\ }
\newcommand{\resp}{respectively\ }
\newcommand{\Subsection}{\subsection}
\providecommand{\nobreakdash}{\nobreak\hbox{-}}
\setlength{\emergencystretch}{2em}
\multlinegap0pt
\usepackage{bm}
\usepackage{bbm}
\usepackage{dsfont}
\usepackage{xspace}
\usepackage{cancel}
\usepackage{mathtools}
\usepackage{amsthm}
\makeatletter
\@ifundefined{theorem}{\newtheorem{theorem}{Theorem}}{}
\makeatother
\newcommand{\HH}{\mathbb{H}} 
\newcommand{\OO}{\mathbb{O}}
\title[Computing the Cousin-Zuckerman Resolution and the Lusztig-Vo]{Computing the Cousin-Zuckerman Resolution and the Lusztig-Vogan Bijection}
\author{Jack A. Cook}
\date{Source: arXiv:2604.19554v1 (CC BY 4.0)}
\begin{document}
\maketitle

\noindent\emph{Source and licence.} Computing the Cousin-Zuckerman Resolution and the Lusztig-Vogan Bijection. Version arXiv:2604.19554v1 (CC BY 4.0), distributed under the Creative Commons Attribution 4.0 International licence, as recorded in the arXiv licence metadata for this exact version and shown on the abstract page https://arxiv.org/abs/2604.19554v1. Full rights notice: licenses/arxiv-2604.19554-CC-BY-4.0.txt.\par\smallskip

\noindent\textit{Editorial scope.} Counted unit: the Theorem whose source label is \texttt{None} in the article (source0.tex lines 2352--2371, source label \texttt{None}), delivered together with the article's own complete original proof of it (source0.tex lines 2373--2396, one \texttt{proof} environment of 24 lines). No mathematics is written, repaired or re-derived: statement and proof are the article's own text line for line. Required context retained uncounted: none. Every definition, display and label of those lines is retained unchanged. Omitted from this package and not redistributed: 1--2351, 2372--2372, 2397--2564. Nothing inside a retained excerpt is omitted or altered: every retained body is delivered exactly as the article's lines are written, so no omission receipt of any kind accompanies this package. Citations: the retained excerpts carry no citation command at all, so no bibliography entry of the article is transcribed and no reference is redistributed. Reference adaptation: none was needed and none was made; every reference inside the delivered unit resolves inside it, so no unresolved reference or citation is produced. Printed slips kept literal: no printed word, symbol, hypothesis or number of the counted statement or of its proof was altered. Layout and class: the article's own class is \texttt{article} with packages including authblk, amssymb, amsmath, stmaryrd, mathrsfs, hyperref, titling, bbm, dsfont, bbold, fontenc, palatino, mathpazo, marvosym; the wrapper is the lane's pinned \texttt{amsart} editorial wrapper at 10pt on A4, which restates the theorem-like environments the retained text needs and adds the article's own macro definitions that the retained text needs (\texttt{\textbackslash HH}, \texttt{\textbackslash OO}), with extra wrapper packages (bm, bbm, dsfont, xspace, cancel, mathtools). The delivered numbering is the unit's own. Assets: no retained span contains a figure environment, a graphics-inclusion command, a PDF-page inclusion, a table or a TikZ picture; no image file is copied into this package, the package declares no asset dependency and no image file is redistributed with it. Licence: version arXiv:2604.19554v1 of this article is distributed under the Creative Commons Attribution 4.0 International licence, as recorded in the arXiv licence metadata for this exact version and shown on the abstract page https://arxiv.org/abs/2604.19554v1. A full rights notice is delivered at licenses/arxiv-2604.19554-CC-BY-4.0.txt. Characters: every retained excerpt is pure ASCII, so the delivered engine, XeLaTeX with its default Latin Modern text font, reports no missing character.
\begin{theorem}
    For $ GL(3,\HH) $, the following assignments are the bijection in $
    (1): $
    \begin{align*}
        (\OO_0,(a,b,c))&\mapsto (a+4,b+2,c) \\
        (\OO_{mid}, (n_1,n_2))& \mapsto {%
        \begin{cases}
            \left(n_2+2,\left \lceil \frac{n_1}{2} \right\rceil, \left
            \lceil \frac{n_1-1}{2} \right\rceil\right) & n_1\leq 2n_2+1
            \\
            \; \\
            \left(\left \lceil \frac{n_1}{2} \right\rceil+1, \left
            \lceil \frac{n_1-1}{2} \right\rceil+1, n_2\right) & n_1>2n_2+1
        \end{cases}
        } \\
        (\OO_{prin},n)&\mapsto \left( \left \lceil \frac{n}{3} \right\rceil,
        \left \lceil \frac{n-1}{3} \right\rceil, \left \lceil \frac{n-2}
        {3} \right\rceil \right).
    \end{align*}
\end{theorem}

\begin{proof}
    The final piece to check is that this map is indeed a bijection.  We
    can classify all of the $ K $-dominant $ T $-weights in the
    following way.  Let $ (x,y,z) $ the weight.  The options for the
    differences $ x-y $ and $ y-z $ are first split into two cases:
    both $ x-y $ and $ y-z $ are $ \geq 2 $ or not.  In the affirmative,
    the only option is to be a weight attached to the zero orbit.  The
    map on the zero orbit is injective as the inputs there must satisfy $
    a\geq b\geq c\geq 0. $

    Suppose then that the $ x-y\geq 2. $ Necessarily the other
    difference must be $ \leq 1 $ and $ x\geq 2 $ by the dominance
    condition.  This immediately rules out the pairs coming from the
    principal orbit as the differences in weights there are identically $
    0 $ or $ 1. $ Thus, this weight must live on the middle orbit.  In
    particular, it must be mapped to by a pair $ (n_1,n_2) $ with $ n_1\leq
    2n_2+1. $ However, in this regime, the map is injective clearly.  A
    similar argument reduces the case of $ x-y\leq 1 $ and $ y-z\geq 2 $
    to the other regime for the middle orbit.  This is also injective
    and thus the entire map is.

    Surjectivity follows at once from the consideration of these cases.

\end{proof}

\end{document}
